\chapter{Survey: [Nom del lloc] -- [Data]}
\label{chap:survey-template}

% Duplica aquest fitxer per a cada nou survey, anomena'l per exemple
% survey_2026-08-04_nom-lloc.tex, i afegeix un \input{} a main.tex.

\section{Objectiu}
\textit{[Pendent de redactar]}

\section{Ubicació}
% Coordenades, mapa de la zona, condicions del terreny.
\textit{[Pendent de redactar]}

\section{Condicions meteorològiques}
\begin{table}[H]
    \centering
    \caption{Condicions durant el vol}
    \label{tab:condicions-survey}
    \begin{tabular}{lc}
        \toprule
        Paràmetre & Valor \\
        \midrule
        Data i hora &  \\
        Temperatura &  \\
        Vent &  \\
        Cobertura núvols &  \\
        \bottomrule
    \end{tabular}
\end{table}

\section{Pla de vol}
% Paràmetres específics d'aquest survey (àrea, altura, solapaments).
\textit{[Pendent de redactar]}

\section{Punts de control (GCP)}
\textit{[Pendent de redactar]}

\section{Execució del vol}
\textit{[Pendent de redactar]}

\section{Processament de dades}
\textit{[Pendent de redactar]}

\section{Resultats}
% Ortofoto, model 3D, DEM/DSM, mètriques de precisió.
\textit{[Pendent de redactar]}

\begin{figure}[H]
    \centering
    % \includegraphics[width=0.8\textwidth]{img/surveys/exemple.jpg}
    \caption{[Descripció de la imatge/resultat]}
    \label{fig:survey-exemple}
\end{figure}

\section{Incidències}
\textit{[Pendent de redactar]}

\section{Conclusions}
\textit{[Pendent de redactar]}
